%% LyX 2.4.4 created this file.  For more info, see https://www.lyx.org/.
%% Do not edit unless you really know what you are doing.
\documentclass[oneside,english]{book}
\usepackage[T1]{fontenc}
\usepackage[latin9]{inputenc}
\setcounter{secnumdepth}{3}
\setcounter{tocdepth}{3}
\usepackage{amstext}
\usepackage{amsthm}
\usepackage{amssymb}
\usepackage{geometry}
\geometry{verbose,tmargin=1in,bmargin=1in,lmargin=1.5in,rmargin=1.5in}
\PassOptionsToPackage{normalem}{ulem}
\usepackage{ulem}

\makeatletter
%%%%%%%%%%%%%%%%%%%%%%%%%%%%%% Textclass specific LaTeX commands.
\theoremstyle{definition}
\newtheorem{defn}{\protect\definitionname}
\theoremstyle{plain}
\newtheorem{thm}{\protect\theoremname}
\ifx\proof\undefined\
  \newenvironment{proof}[1][\proofname]{\par
    \normalfont\topsep6\p@\@plus6\p@\relax
    \trivlist
    \itemindent\parindent
    \item[\hskip\labelsep
          \scshape
      #1]\ignorespaces
  }{%
    \endtrivlist\@endpefalse
  }
  \providecommand{\proofname}{Proof}
\fi
\theoremstyle{plain}
\newtheorem{cor}{\protect\corollaryname}

%%%%%%%%%%%%%%%%%%%%%%%%%%%%%% User specified LaTeX commands.
\usepackage{hyperref}
\usepackage[usenames,dvipsnames]{xcolor} %adds additional color options
\hypersetup{colorlinks=true, urlcolor=Blue}%Blue

\usepackage{multicol}

\usepackage{tikz}
\usepackage{tikz-3dplot}
\usetikzlibrary{calc,arrows,shapes,backgrounds}

\usetikzlibrary{patterns}
\usetikzlibrary{plotmarks}
\usepackage{pgfplots}
\pgfplotsset{compat=newest}

\usepackage{calc}
\usepackage[nomessages]{fp}

\usepackage{datetime}
%\usepackage{subcaption}

%\usepackage{amsthm}
% Added by lyx2lyx
\setlength{\parskip}{\smallskipamount}
\setlength{\parindent}{0pt}

\makeatother

\usepackage{babel}
\providecommand{\corollaryname}{Corollary}
\providecommand{\definitionname}{Definition}
\providecommand{\theoremname}{Theorem}

\begin{document}

\chapter{Ideals\label{chap:Ideals}}

\section{Ideals}
\begin{defn}
A subring $A$ of a ring $R$ is an (two-sided) \emph{ideal }of $R$
if for every $r\in R$ and every $a\in A$ then $ra$ and $ar$ are
in $A.$ 
\end{defn}
Example \ref{3Z is an ideal} is an ideal while example \ref{=00007B0,1=00007D is not an ideal}
is not an ideal.

Of course it will not be necessary to check the property ``$ra$
and $ar$'' when the ring is commutative (we anyway in example \ref{3Z is an ideal}).
However keep in mind that the subring of a non-commutative ring may
still be commutative --in which case the property must be checked.
\begin{thm}
\textbf{\emph{\label{thm:Ideal-test.-A}Ideal test. }}A nonempty subset
$I$ of a ring $R$ is an ideal if and only if:

\begin{enumerate}
\item If $a,b\in I$ then $a-b\in I$.
\item If $r\in R$ and $a\in I$ then $ra\in I$ and $ar\in I$.
\end{enumerate}
\begin{proof}
($\Leftarrow$) Let $I$ have properties $1$ and $2$. $I$ is a
ring by the subring test since property 2 is true for all $r\in I\subseteq R$,
and property 1 is identical to the property 1 of the subring test.
Property 2 also implies the absorbing property of ideals.

($\Rightarrow$) All ideals have property 1 since they are subrings,
and property 2 by definition.
\end{proof}
\label{thm:principal ideals are ideal}Let $R$ be a commutative ring
with unity\footnote{``with unity'' is equivalent to ``with identity''.},
$c\in R$, and $I$ the set of all multiples of $c$ in $R$ ($I=\{cr\colon r\in R\}$)
Then $I$ is an ideal.
\end{thm}
\begin{defn}
The ideal in theorem \ref{thm:principal ideals are ideal} is called
the \emph{principal ideal generated by $c$, denoted $\left\langle c\right\rangle $}\footnote{Also often denoted $(c)$}\emph{.}

A \emph{prime ideal $P$ }of a commutative ring $R$ is a proper ideal\footnote{An ideal $I$ in $R$ is proper if $I\neq R$, i.e., $I$ is a proper
subset of $R$.} of $R$ such that $a,b\in R$ and $ab\in P$ implies $a\in P$ \uline{or}
$b\in P$.
\end{defn}

\begin{defn}
A \emph{maximal ideal $M$ }of a commutative ring $R$ is a proper
ideal such that whenever $I$ is an ideal of $R$ and $M\subseteq I\subseteq R$
then $I=M$ or $I=R$.
\end{defn}

\section{Congruence and ideals}

We are familiar with the idea of congruence. Ideals can be congruent
as well. 
\begin{defn}
Let $I$ be an ideal in a ring $R$ and let $a,b\in R$. Then $a\equiv b\bmod I$
if $a-b\in I$. We say ``$a$ is congruent to $b$ modulo $I$.
\end{defn}
\begin{thm}
\label{thm:congruence-is-equivalence}Let $I$ be an in ideal in the
ring $R$. Then congruence modulo $I$ in $R$ is an equivalence relation.

\label{thm:cong-prop}Let $I$ be an ideal in the ring $R$. If $a\equiv b\bmod I$
and $c\equiv d\bmod I$, then
\end{thm}
\begin{enumerate}
\item $a+c\equiv b+d\bmod I$
\item $ac\equiv bd\bmod I$
\end{enumerate}

\section{Cosets}

Recall that the elements in $\mathbb{Z}$ each fell into $n$ classes
of $\mathbb{Z}_{n}$ denoted $\left[a\right]$ for an $a\in\mathbb{Z}_{6}$.
For $\mathbb{Z}_{6}$, $2\equiv8\bmod6$ so then $8\in\left[2\right]$.
We can easily find all the elements in $\left[2\right]$ by repeatedly
adding 6 (since we are working with modulo $6$). Similarly we could
find all the elements of each class. 
\begin{eqnarray*}
\left[2\right] & = & \left\{ \dots,-2,0,2,8,14,\dots\right\} \\
\left[3\right] & = & \left\{ \dots,-3,0,3,6,15,\dots\right\} \\
\vdots\\
\left[a\right] & = & \left\{ \dots,a-n,a,a+n,a+2n,\dots\right\} \\
\vdots\\
\left[6\right] & = & \left\{ \dots,-6,0,6,12,\dots\right\} 
\end{eqnarray*}
The elements of each class of $\left[a\right]$ can each be expressed
as $a+kn$ for some $k\in\mathbb{Z}$. It thus would make sense to
write $a+6\mathbb{Z}$ instead of $\left[a\right]$. This is in fact
what we will do when expressing the congruence class of elements in
a ring modulo an ideal.
\begin{defn}
Let $I$ be an ideal in $R$ and $a\in R$, then the \emph{congruence
class of }$a$ \emph{modulo $I$ }is the set of all elements of $R$
that are congruent to $a$ modulo $I$, i.e., 
\begin{eqnarray*}
\left\{ b\in R\colon b\equiv a\bmod I\right\}  & = & \left\{ b\in R\colon b-a\in I\right\} \\
 & = & \left\{ b\in R\colon b=a+i\mbox{ for }i\in I\right\} \\
 & = & \left\{ a+i\colon i\in I\right\} 
\end{eqnarray*}
The congruence class of $a$ modulo $I$ will be denoted $a+I$ (instead
of $\left[a\right]$). The congruence class $a+I$ is also called
the (left) \emph{coset of $I$ in $R$.}
\end{defn}

\begin{defn}
\label{def-factor-ring}If $I$ is an ideal in the ring $R$, then
the set of all cosets of $I$ is called the \emph{factor ring}\footnote{Also called the \emph{quotient ring}.}
$R$ modulo $I$ and is denoted $R/I$.
\end{defn}

\begin{thm}
Let $I$ be an ideal in the ring $R$ and let $a,c\in R$. Then $a\equiv c\bmod I$
if and only if $a+I=c+I$

\begin{proof}
$(\Rightarrow)$ Suppose that $a\equiv c\bmod I$. First we will show
that $a+I\subseteq c+I$. Let $b\in a+I$. By definition, $b\equiv a\bmod I$.
Since $a\equiv c\bmod I$ we know that $b\equiv c\bmod I$ by transitivity
(theorem \ref{thm:congruence-is-equivalence}). Thus $a+I\subseteq c+I$.
Similarly we can show that $c+I\subseteq a+I$ since $a\equiv c\bmod I\Rightarrow c\equiv a\bmod I$.

$(\Leftarrow)$ Suppose that $a+I=c+I$. Since $a\equiv a\bmod I$
by reflexivity (theorem \ref{thm:congruence-is-equivalence}), we
have that $a\in a+I$ and thus $a\in c+I$. By definition we then
have that $a\equiv c\bmod I$.
\end{proof}
\end{thm}

\begin{cor}
Let $I$ be an ideal of the ring $R$. Any two cosets of $I$ are
either disjoint or identical.

\begin{proof}
Suppose that $a+I$ and $b+I$ are not disjoint. Then there is an
element $r\in R$ such that $r\in a+I$ and $r\in b+I$. So then $r\equiv a\bmod I$
and $r\equiv b\bmod I$. By symmetry $r\equiv b\bmod I\Rightarrow b\equiv r\bmod I$.
Then by transitivity $a\equiv b\bmod I$. By the theorem above, we
have that $a+I=b+I$. 
\end{proof}
\end{cor}

\end{document}
